\documentclass[11pt]{article}

\usepackage[utf8]{inputenc}
\usepackage[margin=1in]{geometry}
\usepackage{amsmath}
\usepackage{amssymb}
\usepackage{amsthm}
\usepackage{enumitem}

\title{STAT GU4204 --- HW4 Extra Credit (Questions)}
\author{Jonah Aden}
\date{Out: 05/02 --- Submission optional through 05/05 5pm grace period}

\begin{document}

\maketitle

\noindent\textit{Note: This HW is on most recent material covered after the last HW.
Students are responsible for the final exam from the problems in this HW.
Submission is optional and for extra credit, graded on submission and effort.
Not submitting will not negatively affect the course grade.}

\bigskip

\section*{Section 9.5}

\paragraph{Problem 2.}
Suppose that nine observations are selected at random from the normal distribution
with unknown mean $\mu$ and unknown variance $\sigma^2$, and for these nine
observations it is found that $\bar{X}_n = 22$ and
$\sum_{i=1}^{n}(X_i - \bar{X}_n)^2 = 72$.
\begin{enumerate}[label=(\alph*)]
    \item Carry out a test of the following hypotheses at the level of significance $0.05$:
    \[
        H_0\colon \mu \leq 20, \qquad H_1\colon \mu > 20.
    \]
    \item Carry out a test of the following hypotheses at the level of significance $0.05$
    by using the two-sided $t$ test:
    \[
        H_0\colon \mu = 20, \qquad H_1\colon \mu \neq 20.
    \]
    \item From the data, construct the observed confidence interval for $\mu$ with
    confidence coefficient $0.95$.
\end{enumerate}

\paragraph{Problem 4.}
Suppose that a random sample of eight observations $X_1, \dots, X_8$ is taken from
the normal distribution with unknown mean $\mu$ and unknown variance $\sigma^2$,
and it is desired to test the following hypotheses:
\[
    H_0\colon \mu = 0, \qquad H_1\colon \mu \neq 0.
\]
Suppose also that the sample data are such that $\sum_{i=1}^{8} X_i = -11.2$ and
$\sum_{i=1}^{8} X_i^2 = 43.7$. If a symmetric $t$ test is performed at the level of
significance $0.10$ so that each tail of the critical region has probability $0.05$,
should the hypothesis $H_0$ be rejected or not?

\section*{Section 9.6}

\paragraph{Problem 1.}
In Example 9.6.3, we discussed Roman pottery found at two different locations in
Great Britain. There were samples found at other locations as well. One other
location, Island Thorns, had five samples $X_1, \dots, X_n$ with an average
aluminum oxide percentage of $\bar{X} = 18.18$ with
$\sum_{i=1}^{5}(X_i - \bar{X})^2 = 12.61$. Let $Y_1, \dots, Y_5$ be the five sample
measurements from Ashley Rails in Example 9.6.3. Test the null hypothesis that the
mean aluminum oxide percentages at Ashley Rails and Island Thorns are the same
versus the alternative that they are different at level $\alpha_0 = 0.05$.

\section*{Additional Question}

\paragraph{Q1.}
Consider the pdf and the hypothesis test given in the textbook in exercise 9.10: \#4.
The pdf is:
\[
    f(x \mid \theta) =
    \begin{cases}
        2(1-\theta)x + \theta, & 0 \leq x \leq 1, \\
        0, & \text{otherwise},
    \end{cases}
\]
where the value of $\theta$ is unknown ($0 \leq \theta \leq 2$). Given this
hypothesis test rejects $H_0$ when $x > 1/3$, compute the probabilities of Type~I
and Type~II errors.

\end{document}
